\chapter{Конструкторская часть}

\section{Структуры данных}

Для реализации алгоритмов были использованы следующие структуры данных:
\begin{itemize}[label=---]
	\item матрица~---~массив массивов целого типа;
	\item размерности матрицы~---~значения целого типа.
\end{itemize}

\section{Алгоритмы}

Ниже представлены:

\begin{enumerate}[label=\arabic*)]
	\item Алгоритм стандартного умножения матриц~(рисунок \ref{img:standartMul}).
	\item Стандартный алгоритм Винограда~(рисунок \ref{img:standartWinograd}).
	\item Оптимизированный алгоритм Винограда~(рисунок \ref{img:standartWinograd}).
\end{enumerate}


\begin{figure}[H]
	\centering
	\includesvg[inkscapelatex=false, width=0.8\textwidth, height=0.8\textheight]{inc/img/standartMul.svg}
	\caption{Стандартный алгоритм умножения матриц}
	\label{img:standartMul}
\end{figure}

\newpage
\begin{figure}[H]
	\centering
	\includesvg[inkscapelatex=false, width=0.9\textwidth, height=0.9\textheight]{inc/img/standartWinograd.svg}
	\caption{Стандартный алгоритм Винограда}
	\label{img:standartWinograd}
\end{figure}

\newpage
\begin{figure}[H]
	\centering
	\includesvg[inkscapelatex=false, width=0.9\textwidth, height=0.9\textheight]{inc/img/optWinograd.svg}
	\caption{Оптимизированный алгоритм Винограда}
	\label{img:optWinograd}
\end{figure}




